\documentclass[12pt]{article}

\usepackage{fontawesome}
\usepackage{hyperref}
\usepackage{xurl}
\usepackage[tagged, highstructure]{accessibility}
\usepackage{bookmark}
\usepackage{enumitem}

\hypersetup{
    colorlinks=false,
    pdfborder={0 0 0},
    pdftitle={Files and Paths},
    pdfauthor={Adrianna Holden-Gouveia},
    pdfsubject={Introduction to Linux - Lab Assignment},
    pdflang={en-US},
    pdfkeywords={lab assignment, Files and Paths}
}

% Configure PDF bookmarks for navigation
\bookmarksetup{
    numbered,
    open,
}

% Configure list spacing for better accessibility
\setlist{nosep}



\title{Files and Paths}
\author{
        Adrianna Holden-Gouveia \\
        Website: \href{https://aholdengouveia.name}{aholdengouveia.name}\\ 
        \date{\vspace{-5ex}}
        %Email: \href{mailto:admin@aholdengouveia.name}{admin@aholdengouveia.name} \\
%        \faLinkedin{: aholdengouveia} \\
        \faGithub {: aholdengouveia} \\
%        \faTwitter {: aholdengouveia} \\
        }

%S\date{\today}


\begin{document}    

\maketitle

\section*{Accessibility Notice}
This document is also available in HTML format at:

Link: \href{https://aholdengouveia.name/IntroLinux/labs/filesandpaths.html}{\textbf{https://aholdengouveia.name/IntroLinux/labs/filesandpaths.html}}

The HTML version provides enhanced accessibility features including keyboard navigation, screen reader support, responsive design, dark mode support, and high contrast options.


%\begin{abstract}
%This is a template for Linux Administration Lab work
%\end{abstract}
%\tableofcontents

\section*{Objectives:}
\begin{itemize}
    \item Navigate the Linux file system using both absolute and relative paths
    \item Perform basic file management operations, including creating, copying, moving, and deleting files and directories
\end{itemize}
\section*{References}

References, a video, a PowerPoint and some notes are available at my website
Link: \href{https://www.aholdengouveia.name/IntroLinux/filesandpaths.html}{\textbf{https://www.aholdengouveia.name/IntroLinux/filesandpaths.html}}



\subsection*{Complete the following problems:}
Make sure to take a screenshot of each problem solved.  You can use the clear command to make sure the terminal is focused on your current solution
    \begin{enumerate}
        \item Use pwd command to display and record the absolute path of your current working directory.
        \item Create two new directories using mkdir
        \item Use touch command to create five files
        \item Navigate into one of the  directories using cd 
        \item Move one file from your home directory into your current folder using a relative path
        \item Use ls to confirm it moved
        \item Move one file into the folder using its absolute path
        \item Use ls to confirm it moved
        \item Create a duplicate of one file in your folder and name the copy FILENAMEbackup using the cp command and relative paths
        \item Delete one file in your home folder using rm
        \item Perform a long listing ls -l of the home directory to view the final file details and permissions
    \end{enumerate}



\section*{Deliverables}
A document containing the exact commands used for each numbered step, include clearly labelled screenshots of solution showing the successful execution of these commands as proof of testing.
\end{document}